%/***************************************************************************
% *   ex10.tex                                 Version 20240616.122956      *
% *                                                                         *
% *                                                                         *
% *   Copyright (C) 2024         by Gabriel Bulle De Barros                 *
% ***************************************************************************
% *   Permission is granted to copy, distribute and/or modify this document *
% *   under the terms of the GNU Free Documentation License, Version 1.3    *
% *   or any later version published by the Free Software Foundation;       *
% *   with no Invariant Sections, no Front-Cover Texts, and                 *
% *   no Back-Cover Texts.                                                  *
% *   A copy of the license is included in the section entitled             *
% *   "GNU Free Documentation License", or available online at:             *
% *   https://www.gnu.org/copyleft/fdl.html                                 *
% *                                                                         *
% *   This document is distributed in the hope that it will be useful,      *
% *   but WITHOUT ANY WARRANTY; without even the implied warranty of        *
% *   MERCHANTABILITY or FITNESS FOR A PARTICULAR PURPOSE.  See the         *
% *   GNU Free Documentation License for more details.                      *
% *                                                                         *
% *   You should have received a copy of the GNU Free Documentation License *
% *   along with this document; if not, write to the                        *
% *   Free Software Foundation, Inc.,                                       *
% *   59 Temple Place - Suite 330, Boston, MA  02111-1307, USA.             *
% ***************************************************************************
% *   To contact the author, please write to:                               *
% *   Gabriel Bulle De Barros                                               *
% *   Email: gbb1@poli.br                                                   *
% *   Webpage: http://beco.poli.br/gbb1@poli.br                             *
% *   Phone: 81991396798                                                    *
% ***************************************************************************/

\documentclass[a4paper,12pt]{article}

%%%%%%%%%%%%%%%%%%%%%%%%%%%%%%%%%%%%%%%%%%%%%%%%%%%%%%%%%%%%%%%%%%%%%%%%%%%%%%%%
% pacotes
\usepackage[brazil]{babel}
\usepackage[utf8]{inputenc}
\usepackage[T1]{fontenc}
\usepackage{graphicx}
\usepackage{geometry}
\usepackage{listings}
\usepackage{xcolor}
\usepackage{courier}

%%%%%%%%%%%%%%%%%%%%%%%%%%%%%%%%%%%%%%%%%%%%%%%%%%%%%%%%%%%%%%%%%%%%%%%%%%%%%%%%
% comandos

\geometry{a4paper, margin=1in}

%%%%%%%%%%%%%%%%%%%%%%%%%%%%%%%%%%%%%%%%%%%%%%%%%%%%%%%%%%%%%%%%%%%%%%%%%%%%%%%%
% Configurações para o código em C++
\lstset{
    language=C++,
    basicstyle=\small\ttfamily,
    keywordstyle=\bfseries\color{blue},
    commentstyle=\itshape\color{gray},
    stringstyle=\color{green!60!black},
    showstringspaces=false,
    numbers=left,
    numberstyle=\tiny,
    numbersep=5pt,
    frame=single,
    breaklines=true,
    captionpos=b,
    extendedchars=true,
    tabsize=2,
    xleftmargin=0.5cm, % Margem esquerda
    xrightmargin=0.5cm, % Margem direita
}


%%%%%%%%%%%%%%%%%%%%%%%%%%%%%%%%%%%%%%%%%%%%%%%%%%%%%%%%%%%%%%%%%%%%%%%%%%%%%%%%
\title{Manual Técnico da Calculadora Polonesa Reversa}
\author{Arthur Bezerra Melo, Gabriel Bullé de Barros}
\date{16 de junho de 2024}

%%%%%%%%%%%%%%%%%%%%%%%%%%%%%%%%%%%%%%%%%%%%%%%%%%%%%%%%%%%%%%%%%%%%%%%%%%%%%%%%
\begin{document}

\maketitle

\begin{figure}[h]
    \centering
    \includegraphics[width=0.5\linewidth]{logo-poli.png}
    \caption{Universidade de Pernambuco}
    \label{fig:logo-poli.png}
\end{figure}

%%%%%%%%%%%%%%%%%%%%%%%%%%%%%%%%%%%%%%%%%%%%%%%%%%%%%%%%%%%%%%%%%%%%%%%%%%%%%%%%
\section{Introdução e Visão Geral}

Este manual descreve o software da Calculadora Polonesa Reversa implementada em C++. A versão atual é 1.0.0.

%%%%%%%%%%%%%%%%%%%%%%%%%%%%%%%%%%%%%%%%%%%%%%%%%%%%%%%%%%%%%%%%%%%%%%%%%%%%%%%%
\section{Implementação e Funcionamento}

\subsection{Classe \texttt{Stack}}
A classe \texttt{Stack} implementa uma pilha genérica usando um vetor para armazenar os elementos. Ela fornece métodos para empilhar (\texttt{push}), desempilhar (\texttt{pop}), acessar o topo (\texttt{top}), verificar se está vazia (\texttt{isEmpty}) e obter o tamanho (\texttt{size}).

\vspace{20pt}

\begin{minipage}{\linewidth}
\begin{lstlisting}[caption={Implementação da classe \texttt{Stack}}, label={lst:stack}]
template <typename T>
class Stack
{
    private:
        std::vector<T> elements;

    public:
        void push(const T &element)
        {
            elements.push_back(element);
        }
         void pop()
        {
            if(elements.empty())
                throw std::out_of_range("Stack<>::pop(): pilha vazia");
            elements.pop_back();
        }

        T top() const
        {
            if(elements.empty())
                throw std::out_of_range("Stack<>::top(): pilha vazia");
            return elements.back();
        }

        bool isEmpty() const
        {
            return elements.empty();
        }

        size_t size() const
        {
            return elements.size();
        }

};
\subsection{Classe \texttt{PolishCalculator}}
A classe \texttt{PolishCalculator} avalia expressões matemáticas na notação polonesa reversa. Ela utiliza uma instância da classe \texttt{Stack} para manipular os operandos e operadores durante a avaliação.

\vspace{20pt}

\begin{minipage}{\linewidth}
\begin{lstlisting}[caption={Implementação da classe \texttt{PolishCalculator}}, label={lst:polishcalculator}]
// Implemento da calculadora polonesa
class PolishCalculator
{
    public:
        double evaluate(const std::string &expression)
        {
            std::istringstream tokens(expression);
            std::string token;
            Stack<double> stack;

            while(tokens >> token)
            {
                if(isOperator(token))
                {
                    double b = stack.top();
                    stack.pop();
                    double a = stack.top();
                    stack.pop();
                    double result = applyOperator(token, a, b);
                    std::cout << "Aplicando operador " << token << " a " << a << " e " << b << " resulta em " << result << std::endl;
                    stack.push(result);
                }
     [...]
};
\end{lstlisting}
\end{minipage}
\subsubsection{Método \texttt{evaluate}}
O método \texttt{evaluate} recebe uma expressão em formato de string e a avalia passo a passo:
\begin{itemize}
    \item Separa os tokens da expressão usando um \texttt{std::istringstream}.
    \item Para cada token:
      \begin{itemize}
        \item Empilha números na pilha.
        \item Quando encontra um operador, desempilha os dois operandos mais recentes, aplica o operador e empilha o resultado de volta.
      \end{itemize}
    \item Ao final, verifica se há exatamente um valor na pilha. Se não, lança uma exceção indicando uma "Expressão inválida".
\end{itemize}
\newpage
\section{Função \texttt{main}}
A função \texttt{main} controla a execução do programa, interagindo com o usuário e utilizando a calculadora polonesa reversa.

\begin{lstlisting}[caption={Implementação da função principal \texttt{Main}}, label={lst:main}]
int main()
{
    PolishCalculator calculator;
    std::string expression;

    std::cout << "Digite uma expressao em notacao polonesa reversa: ";
    std::getline(std::cin, expression); // Leitura da expressao do usuario

    try
    {
        double result = calculator.evaluate(expression);
        std::cout << "Resultado: " << result << std::endl;
    }
    catch(const std::exception &ex)
    {
        std::cerr << "Error: " << ex.what() << std::endl;
    }

    return 0;
}
\end{lstlisting}

\section{Instruções de Compilação, Uso e Resolução de Problemas}

\subsection{Instruções de Compilação}
Para compilar o programa, use o seguinte comando:
\begin{verbatim}
$make ex10.out
\end{verbatim}

\subsection{Instruções de Uso}
Para usar o programa, execute o binário gerado e digite uma expressão em notação polonesa reversa quando solicitado.

Exemplo de entrada:
\begin{verbatim}
Digite uma expressão em notação polonesa reversa: 5 10 * 25 / 5 *
\end{verbatim}
Saída esperada:
\begin{verbatim}
Resultado: 10
\end{verbatim}

\subsection{Resolução de Problemas}
Ao usar a calculadora, esteja atento aos seguintes erros:
\begin{itemize}
    \item \textbf{Erro: Expressão inválida - faltam operandos}: Indica que a expressão fornecida não possui operandos suficientes para aplicar um operador.
    \item \textbf{Erro: Divisão por zero}: Indica que houve uma tentativa de dividir por zero durante a avaliação da expressão.
    \item \textbf{Erro: Operador inválido}: Indica que foi encontrado um operador desconhecido ou inválido na expressão.
\end{itemize}

Para suporte adicional, entre em contato com [gbb1@poli.br, abm4@poli.br].

%%%%%%%%%%%%%%%%%%%%%%%%%%%%%%%%%%%%%%%%%%%%%%%%%%%%%%%%%%%%%%%%%%%%%%%%%%%%%%%%
% Bibliografia
%\bibliography{bibliografia}

%%%%%%%%%%%%%%%%%%%%%%%%%%%%%%%%%%%%%%%%%%%%%%%%%%%%%%%%%%%%%%%%%%%%%%%%%%%%%%%%
% Índice Analítico de Termos
%\index{algoritmo!de retropropagação|see{rede, de retropropagação}}
%\printindex


\vfill

Typed with \LaTeX.

\end{document}

%/* ----------------------------------------------------------------------- */
%/* vi: set ai et ts=4 sw=4 tw=0 wm=0 fo=croql : PL config for Vim modeline */
%/* Template by Dr. Beco <rcb at beco dot cc>       Version 20150620.224740 */


